\begin{center}\begin{tikzpicture}[font=\small,>=stealth,line join=round]
\draw[gray!40,rounded corners=3pt] (0,0) rectangle (14.2,8.3);\node at (7.1,7.700000000000001) {固定指数2，改变系数会改变每一点的高度};
\begin{scope}[shift={(4.1,2.11)},x=1.6666666666666665cm,y=0.9199999999999999cm]\draw[->,gray] (-1.8,0)--(1.8,0) node[right]{$x$};\draw[->,gray] (0,-0.5)--(0,4.5) node[above]{$y$};\draw[blue,thick,solid,domain=-1.35:1.35,samples=180] plot(\x,{(\x)^2});\draw[teal,thick,dashed,domain=-1.35:1.35,samples=180] plot(\x,{2*(\x)^2});\draw[violet,thick,densely dotted,domain=-1.35:1.35,samples=180] plot(\x,{.5*(\x)^2});\fill[orange!85!black] (1,1) circle(2.2pt);\fill[orange!85!black] (1,2) circle(2.2pt);\fill[orange!85!black] (1,0.5) circle(2.2pt);\draw[gray,dashed] (1,0)--(1,3.9);\node[below] at (1,0){$1$};\end{scope}\node[anchor=west,align=left,text width=5.0cm] at (8.4,4.0) {蓝：$x^2$\\[9pt]青：$2x^2$\\[9pt]紫：$\frac12x^2$\\[9pt]$x=1$处高度为$1,2,\frac12$\\[9pt]严格定义只取系数1};
\node[text width=12.6cm,align=center] at (7.1,.65) {同一指数也能有不同高度；固定系数1，才是单独研究指数的基本幂函数。};
\end{tikzpicture}\end{center}
